\begin{table}[htbp]
\centering
\caption{Full scenario policy ladder across 18 states per corridor--objective pair}\label{tab:scenario_policy_recovery}
\begingroup
\small
\setlength{\tabcolsep}{0.55em}
\renewcommand{\arraystretch}{1.08}
\begin{tabular}{@{}llrrrrrrr@{}}
\toprule
Corridor & Objective & P0 & P1 & P2 & P3 & P3 recoveries & P4 & P4$-$P3 \\
\midrule
Hangzhou--Nanjing & Profit & 2 & 3 & 4 & 5 & 3 & 5 & 0 \\
Hangzhou--Nanjing & Welfare change & 2 & 3 & 4 & 5 & 3 & 5 & 0 \\
Wuhan--Changsha & Profit & 3 & 4 & 9 & 13 & 10 & 18 & 5 \\
Wuhan--Changsha & Welfare change & 3 & 4 & 9 & 14 & 11 & 18 & 4 \\
Jiangxi & Profit & 5 & 9 & 15 & 17 & 12 & 18 & 1 \\
Jiangxi & Welfare change & 5 & 9 & 15 & 18 & 13 & 18 & 0 \\
\bottomrule
\end{tabular}
\parbox{0.96\textwidth}{\footnotesize\textit{Notes.} Each policy count is the number of feasible solutions among 18 income--electrification states. A P3 recovery has infeasible P0 and feasible P3; P4$-$P3 is the number of additional K-dependent feasible states.}
\endgroup
\end{table}
